\section{Architecture and Objective Details}
\label{app:architecture_details}

This appendix specifies the implementation-level equations omitted from the
main text. All distances supplied to neural networks are represented by their
numerical values in Angstrom.

\subsection{Node initialization and edge features}
\label{app:edge_features}

\paragraph{Radial basis.}
For an edge distance \(d\), the docking trunk uses 32 Gaussian radial basis
functions
\begin{equation}
    \rho_m(d)
    =
    \exp\left[-\frac{(d-\mu_m)^2}{2\Delta^2}\right],
    \qquad
    \mu_m\in[0,16]\,\angstrom,\quad
    \Delta=\frac{16}{32}\,\angstrom.
    \label{eq:rbf}
\end{equation}

\paragraph{Conditioning and initialization.}
The conditioning vector in Eq.~\eqref{eq:time_sigma_embedding} uses a
32-dimensional sinusoidal embedding. The final linear map of the prior-scale
branch is initialized to zero. Let \(\mathbf{a}_i\) denote raw node
attributes. Initial scalar states are
\begin{align}
    \mathbf{s}^{0}_i
    &=
    \operatorname{MLP}_{\mathrm{node}}(\mathbf{a}_i),
    \qquad i\notin\mathcal{V}_{\mathrm{frag}},\\
    \mathbf{s}^{0}_k
    &=
    \operatorname{MLP}_{\mathrm{frag}}
    \left(
      \left[
      \operatorname{MLP}_{\mathrm{node}}(\mathbf{a}_k),
      \operatorname{Emb}(|\mathcal{A}_k|),
      \boldsymbol{\eta}_{b(k)}
      \right]
    \right),
    \qquad k\in\mathcal{V}_{\mathrm{frag}}.
    \label{eq:scalar_initialization}
\end{align}
With graph center
\(\mathbf{c}_b=|\mathcal{V}_b|^{-1}\sum_{i\in\mathcal{V}_b}\mathbf{x}_i\),
the \(c\)-th initial \(1o\) channel is
\begin{equation}
    \mathbf{h}^{0,1o}_{i,c}
    =
    G_c(\mathbf{s}^{0}_i)(\mathbf{x}_i-\mathbf{c}_{b(i)})
    +
    \mathbb{1}_{i\in\mathcal{V}_{\mathrm{frag}}}
    \sum_{q=1}^{3}W^{R}_{cq}\mathbf{R}_{i,t}^{(:,q)}.
    \label{eq:vector_initialization}
\end{equation}
Here \(G_c\) ends in \(\tanh\), \(\mathbf{R}^{(:,q)}\) is a rotation-matrix
column, and \(W^R\) is zero-initialized.

\paragraph{Complete edge scalar.}
Let \(\bar{\mathbf{s}}_i^{(\ell)}\in\R^{384}\) be the scalar block after
pre-normalization. The reference-distance feature is
\begin{equation}
    \boldsymbol{\delta}_{ji}
    =
    \left[
      d_{ji},\;
      d_{ji}-d^{\mathrm{ref}}_{ji},\;
      q_{ji}\log
      \frac{d_{ji}}{d^{\mathrm{ref}}_{ji}+\epsilon},\;
      q_{ji}
    \right],
    \qquad
    q_{ji}=\mathbb{1}[d^{\mathrm{ref}}_{ji}>0].
    \label{eq:distance_evolution}
\end{equation}
For a source assigned to fragment \(k\), the local-frame feature is
\(\mathbf{R}_{k,t}^{\top}\mathbf{r}_{ji}\); it is zero otherwise. The
six-dimensional pair vector \(\boldsymbol{\chi}_{ji}\) represents
hydrogen-bond compatibility, salt-bridge compatibility, equal-charge clash,
hydrophobic contact, halogen--acceptor contact, and metal coordination. The
full edge descriptor is
\begin{align}
    \mathbf{e}^{(\ell)}_{ji}
    =\big[
      &\boldsymbol{\rho}(d_{ji}),
      \operatorname{Emb}(\tau_{ji}),
      \operatorname{Emb}(\mathrm{bond}_{ji}),
      W_d\boldsymbol{\delta}_{ji},
      \operatorname{Emb}(\mathrm{hop}_{ji}),\nonumber\\
      &W_R\mathbf{R}_{k,t}^{\top}\mathbf{r}_{ji},
      W_c\boldsymbol{\chi}_{ji},
      \bar{\mathbf{s}}_j^{(\ell)},
      \bar{\mathbf{s}}_i^{(\ell)},
      \boldsymbol{\eta}_{b(i)}
    \big]\in\R^{1020}.
    \label{eq:edge_scalar}
\end{align}
The bond embedding concatenates bond order, conjugation, ring, and
stereochemistry embeddings; unavailable attributes use explicit sentinel
categories.

\subsection{Equivariant normalization and gated aggregation}

\paragraph{Normalization and activation.}
For an \((l,p)\) block
\(\mathbf{h}^{lp}\in\R^{C_{lp}\times(2l+1)}\), equivariant RMS
normalization is
\begin{align}
    \operatorname{rms}_{lp}(\mathbf{h})
    &=
    \left(
      \frac{1}{C_{lp}}
      \sum_{c=1}^{C_{lp}}\|\mathbf{h}^{lp}_c\|_2^2+\epsilon
    \right)^{1/2},\\
    \operatorname{RMSNorm}(\mathbf{h})^{lp}_c
    &=
    \gamma^{lp}_c
    \frac{\mathbf{h}^{lp}_c}{\operatorname{rms}_{lp}(\mathbf{h})}.
    \label{eq:irrep_rmsnorm}
\end{align}
\(0e\) scalars receive SiLU. Non-scalar channels receive invariant norm gates,
\begin{equation}
    \operatorname{EqAct}(\mathbf{h})^{lp}_c
    =
    \mathbf{h}^{lp}_c\,
    \operatorname{sigmoid}
    \left(
      A^{lp}_c(
      \{\|\mathbf{h}^{l'p'}_{c'}\|_2\}_{l'>0,c'})
    \right),
    \qquad l>0.
    \label{eq:equivariant_activation}
\end{equation}
Dropout samples one mask per irrep channel and broadcasts it over the
\(2l+1\) representation components.

\paragraph{Dual radial scales.}
The input and output channel scales in Eq.~\eqref{eq:tp_message} are
\begin{equation}
    \mathbf{u}_{ji}^{(\ell)}
    =
    \operatorname{SiLU}
    (W_r^{(\ell)}\mathbf{e}_{ji}^{(\ell)}+\mathbf{b}_r^{(\ell)}),
    \qquad
    \boldsymbol{\alpha}_{ji}^{\mathrm{in/out}}
    =
    \mathbf{1}
    +W_{\mathrm{in/out}}^{(\ell)}
    \mathbf{u}_{ji}^{(\ell)}.
    \label{eq:dual_radial}
\end{equation}
The scales are broadcast across all \(2l+1\) components of a channel.

\paragraph{Edge gates and aggregation.}
For edge type \(\tau_{ji}\), the gate network returns a scalar component and
one component per non-scalar channel:
\begin{align}
    [a_{ji}^{0},\mathbf{a}_{ji}^{\mathrm{ns}}]
    &=
    \operatorname{MLP}_{g}^{(\ell)}
    (\mathbf{e}_{ji}^{(\ell)}),\\
    g_{ji}^{0}
    &=
    \operatorname{sigmoid}(a_{ji}^{0})
    \exp\left(-\frac{d_{ji}}{\zeta_{\tau_{ji}}^{(\ell)}}\right),
    \qquad
    \zeta_{\tau,\mathrm{init}}^{(\ell)}=8\,\angstrom,\\
    g_{ji,c}^{lp}
    &=
    g_{ji}^{0}\operatorname{sigmoid}(a_{ji,c}^{lp}),
    \qquad l>0.
    \label{eq:edge_gates}
\end{align}
Scalar and non-scalar aggregates are respectively
\begin{align}
    \mathbf{q}_{i,c}^{0e}
    &=
    \frac{
      \sum_{j\in\mathcal{N}(i)}
      g_{ji}^{0}\mathbf{m}_{ji,c}^{0e}
    }{
      \sum_{j\in\mathcal{N}(i)}g_{ji}^{0}+\epsilon
    },
    \label{eq:scalar_aggregation}\\
    \widetilde{\mathbf{q}}_{i,c}^{lp}
    &=
    \frac{
      \sum_{j\in\mathcal{N}(i)}
      g_{ji,c}^{lp}\mathbf{m}_{ji,c}^{lp}
    }{
      \sum_{j\in\mathcal{N}(i)}g_{ji,c}^{lp}+\epsilon
    }.
    \label{eq:nonscalar_aggregation}
\end{align}
Channels with the same \(l\) are split into direction and magnitude. If
\(\boldsymbol{\nu}_{i,l}\) concatenates
\(\|\widetilde{\mathbf{q}}_{i,c}^{lp}\|_2\) across both parities, then
\begin{equation}
    \mathbf{q}_{i,c}^{lp}
    =
    \frac{\widetilde{\mathbf{q}}_{i,c}^{lp}}
         {\|\widetilde{\mathbf{q}}_{i,c}^{lp}\|_2+\epsilon}
    \left[
      \operatorname{SiLU}
      (W_l^{(\ell)}\boldsymbol{\nu}_{i,l}+\mathbf{b}_l^{(\ell)})
    \right]_c.
    \label{eq:nonscalar_rescale}
\end{equation}

\paragraph{Adaptive normalization.}
Let
\(\mathbf{u}_i^{(\ell)}
=\mathbf{h}_i^{(\ell)}
+\operatorname{EqBlock}^{(\ell)}(\mathbf{q}_i^{(\ell)})\) and
\(\bar{\mathbf{u}}=\operatorname{RMSNorm}(\mathbf{u})\). The expanded form of
Eq.~\eqref{eq:compact_layer_update} is
\begin{align}
    \mathbf{h}_{i}^{(\ell+1),0e}
    &=
    \bar{\mathbf{u}}_{i}^{0e}
    \odot(\mathbf{1}+\boldsymbol{\gamma}_{i}^{0e})
    +\boldsymbol{\beta}_{i}^{0e},\\
    \mathbf{h}_{i,c}^{(\ell+1),lp}
    &=
    \bar{\mathbf{u}}_{i,c}^{lp}
    \left[1+0.1\tanh(\gamma_{i,c}^{lp})\right],
    \qquad l>0,\\
    [\boldsymbol{\gamma}_{i}^{0e},
      \boldsymbol{\beta}_{i}^{0e},
      \boldsymbol{\gamma}_{i}^{\mathrm{ns}}]
    &=
    W_{\eta}^{(\ell)}\boldsymbol{\eta}_{b(i)}
    +\mathbf{b}_{\eta}^{(\ell)}.
    \label{eq:equivariant_adaln}
\end{align}
The conditioning projection is zero-initialized.

\subsection{Newton--Euler observable subspace}
\label{app:newton_euler_details}

For
\(\mathbf{I}_k=\mathbf{V}_k
\operatorname{diag}(\boldsymbol{\lambda}_k)\mathbf{V}_k^\top\), define
\begin{align}
    m_{ka}
    &=
    \mathbb{1}[|\mathcal{A}_k|>1]\,
    \mathbb{1}
    \left[
      \lambda_{ka}>0.01\max_b\lambda_{kb}
    \right],\\
    \mathbf{I}_k^{+}
    &=
    \mathbf{V}_k
    \operatorname{diag}
    \left(
      \frac{m_{ka}}{\max(\lambda_{ka},10^{-6})}
    \right)
    \mathbf{V}_k^\top,\\
    \mathbf{P}_k
    &=
    \mathbf{V}_k\operatorname{diag}(\mathbf{m}_k)\mathbf{V}_k^\top.
    \label{eq:observable_projection}
\end{align}
The eigendecomposition is evaluated on a symmetrized double-precision inertia
matrix. Single-atom fragments bypass the eigensolver.

\subsection{Auxiliary docking losses}
\label{app:loss_details}

The atom-level auxiliary objective is
\begin{align}
    \widehat{\mathbf{u}}_i
    &=
    \widehat{\mathbf{v}}_{g(i)}
    +
    \widehat{\boldsymbol{\omega}}_{g(i)}
    \times\mathbf{r}_i,\\
    \mathbf{u}^{\star}_i
    &=
    \mathbf{v}^{\star}_{g(i)}
    +
    \boldsymbol{\omega}^{\star}_{g(i)}
    \times\mathbf{r}_i,\\
    \mathcal{L}_{\mathrm{atom}}
    &=
    \frac{1}{3N}
    \sum_{i=1}^{N}
    \|\widehat{\mathbf{u}}_i-\mathbf{u}^{\star}_i\|_2^2.
    \label{eq:atom_aux_loss}
\end{align}

For the distance-geometry loss, one predicted step gives
\begin{align}
    \widetilde{\mathbf{T}}_{k,1}
    &=
    \mathbf{T}_{k,t_b}
    +(1-t_b)\widehat{\mathbf{v}}_k,\\
    \widetilde{\mathbf{R}}_{k,1}
    &=
    \operatorname{Exp}
    \left((1-t_b)\widehat{\boldsymbol{\omega}}_k\right)
    \mathbf{R}_{k,t_b},\\
    \widetilde{\mathbf{x}}_{i,1}
    &=
    \widetilde{\mathbf{R}}_{g(i),1}\mathbf{y}_i
    +\widetilde{\mathbf{T}}_{g(i),1}.
    \label{eq:one_step_endpoint}
\end{align}
For inter-fragment atom pairs
\(\mathcal{P}_b=\{(i,j):i<j,\ g(i)\neq g(j)\}\),
\begin{equation}
    \mathcal{L}_{\mathrm{DG}}
    =
    \frac{
      \sum_b t_b^2
      \frac{1}{|\mathcal{P}_b|}
      \sum_{(i,j)\in\mathcal{P}_b}
      \left(
        \|\widetilde{\mathbf{x}}_{i,1}
          -\widetilde{\mathbf{x}}_{j,1}\|_2
        -
        \|\mathbf{x}^{\star}_i-\mathbf{x}^{\star}_j\|_2
      \right)^2
    }{
      \sum_b t_b^2
    }.
    \label{eq:distance_geometry_loss}
\end{equation}

\section{Confidence Model Details}
\label{app:confidence_details}

\subsection{Invariant features and readout}

For candidate \(m\), the confidence hidden state is initialized by
\begin{equation}
    \mathbf{h}^{C,0}_{mi}
    =
    \mathbf{h}^{\mathrm{fresh}}_{mi}
    +
    \alpha\,
    \mathbb{1}
    [i\in\mathcal{V}_{\mathrm{lig\text{-}atom}}
      \cup\mathcal{V}_{\mathrm{frag}}]\,
    \mathbf{h}^{D}_{mi},
    \qquad \alpha_{\mathrm{init}}=0.25.
    \label{eq:confidence_hidden_init}
\end{equation}
The invariantized node descriptor is
\begin{equation}
    \mathbf{z}_{mi}
    =
    \left[
      \mathbf{h}_{mi}^{0e},
      \{\|\mathbf{h}_{mi,c}^{lp}\|_2\}_{l>0,p,c}
    \right].
    \label{eq:confidence_invariants}
\end{equation}

For ligand atom \(i\), let
\(d_{mip}=\|\mathbf{x}_{mi}-\mathbf{x}_p\|_2\),
\(d^{\min}_{mi}=\min_p d_{mip}\), and
\(w_{mip}=\exp[-d_{mip}/(2\,\angstrom)]\). With protein-atom flags
\(\mathbf{q}_p\), the explicit contact descriptor is
\begin{align}
    \mathbf{c}_{mi}
    =\bigg[
      &\frac{d^{\min}_{mi}}{10\,\angstrom},
      \boldsymbol{\rho}_{[0,10]}(d^{\min}_{mi}),
      \left\{
      \log\left(1+\sum_p\mathbb{1}[d_{mip}<r]\right)
      \right\}_{r\in\{2,3.5,5\}\,\angstrom},\nonumber\\
      &\log\left(1+\sum_p w_{mip}\right),
      \frac{\sum_p w_{mip}\mathbf{q}_p}
           {\sum_p w_{mip}+\epsilon}
    \bigg].
    \label{eq:confidence_contact_features}
\end{align}
The atom features and outputs are
\begin{align}
    \mathbf{a}_{mi}
    &=
    \operatorname{MLP}_{a}
    \left(
      \operatorname{LN}[
      \mathbf{z}_{mi},\mathbf{c}_{mi}]
    \right),\\
    [\widehat{\ell}^{\,\mathrm{atom}}_{mi},
      \widehat{o}^{\,\mathrm{atom}}_{mi}]
    &=
    \operatorname{MLP}_{\mathrm{atom}}(\mathbf{a}_{mi}).
    \label{eq:confidence_atom_head}
\end{align}

Contact attention is
\begin{align}
    \mathbf{u}_{mi}
    &=
    \operatorname{MLP}_{c}
    \left(
      \operatorname{LN}[
      \mathbf{a}_{mi},\mathbf{c}_{mi}]
    \right),\\
    \alpha_{mi}
    &=
    \frac{
      \exp(s(\mathbf{u}_{mi}))
    }{
      \sum_{j\in\mathcal{V}_{\mathrm{lig\text{-}atom}}}
      \exp(s(\mathbf{u}_{mj}))
    },\\
    \mathbf{r}^{\mathrm{contact}}_m
    &=
    \left[
      \sum_i\alpha_{mi}\mathbf{u}_{mi},
      \max_i\mathbf{u}_{mi},
      \operatorname{mean}_i\mathbf{a}_{mi},
      \max_i\mathbf{a}_{mi}
    \right].
    \label{eq:contact_attention}
\end{align}
For node type
\(q\in\{\mathrm{lig\text{-}atom},\mathrm{frag},
\mathrm{prot\text{-}atom},\mathrm{prot\text{-}res}\}\), let
\(\mathbf{v}_{mi}=\operatorname{MLP}_g(\operatorname{LN}(\mathbf{z}_{mi}))\).
Then
\begin{equation}
    \mathbf{r}^{\mathrm{global}}_m
    =
    \mathop{\Vert}_{q}
    \left[
      \operatorname{mean}_{i:\tau(i)=q}\mathbf{v}_{mi},
      \max_{i:\tau(i)=q}\mathbf{v}_{mi}
    \right].
    \label{eq:global_type_pooling}
\end{equation}
The pose head is defined in Eq.~\eqref{eq:confidence_pose_head}; its reported
RMSD prediction is
\begin{equation}
    \widehat{r}_m
    =
    \max
    \left\{
      0,\,
      \exp\left(
        \operatorname{clip}
        (\widehat{\ell}^{\,\mathrm{pose}}_m,-2,5)
      \right)-1
    \right\}.
\end{equation}

\subsection{Confidence targets and losses}
\label{app:confidence_loss_details}

For fixed atom correspondence,
\begin{equation}
    d_{mi}^{\star}
    =
    \|\mathbf{x}_{mi}-\mathbf{x}^{\star}_i\|_2,
    \qquad
    r_m^{\star}
    =
    \left(
      \frac{1}{N}\sum_i(d_{mi}^{\star})^2
    \right)^{1/2}.
    \label{eq:confidence_targets}
\end{equation}
With unit-transition Huber loss \(H_1\) and binary cross entropy
\(\operatorname{BCE}\), the pointwise terms are
\begin{align}
    \mathcal{L}_{\mathrm{atom\text{-}reg}}
    &=
    \frac{1}{MN}\sum_{m,i}
    H_1\left(
      \widehat{\ell}^{\,\mathrm{atom}}_{mi},
      \log(1+d_{mi}^{\star})
    \right),\\
    \mathcal{L}_{\mathrm{pose\text{-}reg}}
    &=
    \frac{1}{M}\sum_m
    H_1\left(
      \widehat{\ell}^{\,\mathrm{pose}}_m,
      \log(1+r_m^{\star})
    \right),\\
    \mathcal{L}_{\mathrm{atom\text{-}BCE}}
    &=
    \frac{1}{MN}\sum_{m,i}
    \operatorname{BCE}
    \left(
      \widehat{o}^{\,\mathrm{atom}}_{mi},
      \mathbb{1}[d_{mi}^{\star}<2\,\angstrom]
    \right),\\
    \mathcal{L}_{\mathrm{pose\text{-}BCE}}
    &=
    \frac{1}{M}\sum_m
    \operatorname{BCE}
    \left(
      \widehat{o}^{\,\mathrm{pose}}_m,
      \mathbb{1}[r_m^{\star}<2\,\angstrom]
    \right).
    \label{eq:confidence_pointwise_losses}
\end{align}

For
\(\mathcal{Q}=\{(m,n):
\log(1+r_n^{\star})-\log(1+r_m^{\star})>0.3\}\), the pairwise term is
\begin{equation}
    \mathcal{L}_{\mathrm{rank}}
    =
    \frac{1}{|\mathcal{Q}|}
    \sum_{(m,n)\in\mathcal{Q}}
    \max
    \left(
      0,\,
      0.05-
      (\widehat{\ell}^{\,\mathrm{pose}}_n
       -\widehat{\ell}^{\,\mathrm{pose}}_m)
    \right),
    \label{eq:confidence_pairwise}
\end{equation}
with zero contribution when \(\mathcal{Q}\) is empty. The success-listwise
target and loss are
\begin{align}
    q_m
    &=
    \frac{
      \operatorname{sigmoid}
      ((2-r_m^{\star})/0.2)
    }{
      \sum_n
      \operatorname{sigmoid}
      ((2-r_n^{\star})/0.2)
    },\\
    \mathcal{L}_{\mathrm{success\text{-}listwise}}
    &=
    -\sum_m q_m
    \log
    \frac{
      \exp(\widehat{o}^{\,\mathrm{pose}}_m)
    }{
      \sum_n\exp(\widehat{o}^{\,\mathrm{pose}}_n)
    }.
    \label{eq:confidence_listwise}
\end{align}

\section{Additional Reproducibility Material}

The active benchmark and sensitivity studies use the promoted geometry
checkpoint at step 100,000 (SHA-256 prefix \texttt{6932fb3ba6eb}) and
confidence checkpoint at step 42,500 (prefix \texttt{e31fde6f3512}). The
training-configuration hash begins with \texttt{39aa62e4a48e}. Full
checkpoint, configuration, and dataset-manifest hashes are retained in the
machine-readable benchmark summaries rather than abbreviated during
execution.

The active runs used PyTorch 2.10.0, CUDA 13.0, and RTX 6000 Ada or H100 PCIe
accelerators. Every final benchmark and pocket-sensitivity condition contains
all frozen IDs with zero unresolved failures. Exact wall-clock and peak-memory
measurements remain to be collected under one dedicated hardware
configuration.

Machine-readable summaries retain the first, Vina+DG, pure-confidence,
historical-composite, and oracle selectors; RMSD thresholds at 1, 2, 3, and
\(5\,\angstrom\); heavy-atom, fragment-count, rotatable-bond, and cofactor
slices; failure and rescue records; runtime versions; and output paths. The
public archival location for these artifacts will be added with the final
code and weight release.
